% Chapter 2 continuation: source PDF pages 52--59 / printed pages 40--47.
\hypertarget{chapter-02-page-052}{}
对实直线上的函数 $f$, 定义单侧 Hardy--Littlewood 极大函数
\[
 M^+f(t)=\sup_{h>0}\frac1h\int_t^{t+h}|f(s)|\,ds,
 \qquad
 M^-f(t)=\sup_{h>0}\frac1h\int_{t-h}^{t}|f(s)|\,ds.
\]
Hardy 与 Littlewood 所定义的极大函数, 对 $\mathbb R^+$ 上的函数对应于 $M^-$(定义中 $0<h<t$). 当 $|f|$ 递减时, 这一极大函数与作用于 $|f|$ 的 Hardy 算子一致.

若 $f$ 是 $\mathbb R^n$ 上的可测函数, 可在 $(0,\infty)$ 上定义一个递减函数 $f^*$, 称为 $f$ 的递减重排, 它与 $f$ 有相同的分布函数:
\[
 f^*(t)=\inf\{\lambda:a_f(\lambda)\le t\}.
\]
由于 $f$ 与 $f^*$ 有相同的分布函数, 由 \eqref{eq:ch2-distribution-layer-cake}, 它们的 $L^p$ 范数相等; 任何只依赖分布函数的量也相等(比较下文第~\ref{subsec:ch2-lorentz-spaces}~节). Hardy 算子作用于 $f^*$ 的结果通常记作 $f^{**}$.

Hardy 与 Littlewood 证明, 对 $\mathbb R^+$ 上的函数,
\begin{equation}
\label{eq:ch2-one-sided-rearrangement}
 |\{x:M^-f(x)>\lambda\}|\le |\{x:f^{**}(x)>\lambda\}|.
\end{equation}
因此 $M^-$ 的弱 $(1,1)$ 与强 $(p,p)$ 不等式可由 $T$ 的相应不等式推出.

F. Riesz 利用他的“上升太阳引理”给出了 $M^+$ 弱 $(1,1)$ 不等式的一个优美证明(\emph{Sur un théorème du maximum de MM. Hardy et Littlewood}, J. London Math. Soc. 7 (1932), 10--13). 给定 $\mathbb R$ 上的函数 $F$, 若存在 $y>x$ 使 $F(y)>F(x)$, 则称 $x$ 为 $F$ 的阴影点. $F$ 的阴影点集记作 $S(F)$.

\begin{Lemma}
\label{lem:ch2-rising-sun}
设 $F$ 为连续函数, 且
\[
 \lim_{x\to+\infty}F(x)=-\infty,
 \qquad
 \lim_{x\to-\infty}F(x)=+\infty.
\]
则 $S(F)$ 为开集, 并可写成有限开区间的两两不交并 $\bigcup_j(a_j,b_j)$, 且 $F(a_j)=F(b_j)$.
\end{Lemma}

给定 $f\in L^1(\mathbb R)$ 与 $\lambda>0$, 定义 $F(x)=\int_0^x|f(t)|\,dt-\lambda x$. 则 $E(\lambda)=\{x:M^+f(x)>\lambda\}=S(F)$. 利用引理~\ref{lem:ch2-rising-sun} 可推出
\[
 E(\lambda)=\frac1\lambda\int_{E(\lambda)}|f(x)|\,dx,
\]
这与 Hardy 算子的 \eqref{eq:ch2-hardy-level-set} 类似. 强 $(p,p)$ 不等式可像以前一样推出, 常数为 $p'$(且为最佳常数). 证明细节及引理~\ref{lem:ch2-rising-sun} 的证明留给读者.

\hypertarget{chapter-02-page-053}{}
对 $\mathbb R^n$ 上的极大算子, 有与 \eqref{eq:ch2-one-sided-rearrangement} 类似的不等式. 事实上, 存在正常数 $c_n,C_n$, 使
\[
 c_n(Mf)^*(t)\le f^{**}(t)\le C_n(Mf)^*(t),
 \qquad t\in\mathbb R^+.
\]
左侧不等式归功于 F. Riesz; 右侧不等式在 $n=1$ 时归功于 C. Herz, 在 $n>1$ 时归功于 C. Bennett 与 R. Sharpley. 关于证明和更多参考资料, 见后二位作者上述著作的第 3 章.

\subsubsection{Lorentz 空间 $L^{p,q}$}
\label{subsec:ch2-lorentz-spaces}
设 $(X,\mu)$ 为测度空间. $L^{p,q}(X)$ 表示满足下列条件的可测函数 $f$ 所成的空间:
\[
 \|f\|_{p,q}=\left(\frac qp\int_0^\infty[t^{1/p}f^*(t)]^q\frac{dt}{t}\right)^{1/q}<\infty,
 \quad 1\le p<\infty,\ 1\le q<\infty,
\]
以及
\[
 \|f\|_{p,\infty}=\sup_{t>0}t^{1/p}f^*(t)<\infty,
 \quad 1\le p\le\infty.
\]
当 $p=q$ 时,
\[
 \|f\|_{p,p}=\|f^*\|_p=\|f\|_p,
\]
从而恢复 $L^p$. 一般而言, $\|\cdot\|_{p,q}$ 并非范数, 因为三角不等式仅在 $1\le q\le p<\infty$ 或 $p=q=\infty$ 时成立. 但当 $1<p\le\infty$ 且 $1\le q\le\infty$ 时, 若在定义中以 $f^{**}$ 代替 $f^*$, 所得量与 $\|f\|_{p,q}$ 等价并定义一个范数.

若 $q_1\le q_2$, 则 $\|f\|_{p,q_2}\le\|f\|_{p,q_1}$, 因而 $L^{p,q_1}\subset L^{p,q_2}$. 算子 $T$ 为弱 $(p,p)$ 型, 当且仅当
\[
 \|Tf\|_{p,\infty}\le C\|f\|_p.
\]
Marcinkiewicz 插值定理可推广到这些空间. 例如, 它给出比推论~\ref{cor:ch1-hausdorff-young} 更强的 Hausdorff--Young 不等式: 若 $f\in L^p(\mathbb R^n)$, $1<p\le2$, 则 $\widehat f\in L^{p',p}$, 且存在常数 $\bar B_p$, 使
\[
 \|\widehat f\|_{p',p}\le \bar B_p\|f\|_p.
\]
关于递减重排和 $L^{p,q}$ 空间的更多资料, 见 Stein 与 Weiss \MBUnresolvedReference{citation}{[18, Chapter 5]} 及 Bennett 与 Sharpley 的上述著作. 关于插值定理, 特别是非常适用于 Lorentz 空间的所谓实插值法, 也可见后一部著作.

\hypertarget{chapter-02-page-054}{}
这些空间由 G. G. Lorentz 在两篇论文中引入: \emph{Some new functional spaces}, Ann. of Math. 51 (1950), 37--55, 以及 \emph{On the theory of spaces $\Lambda$}, Pacific J. Math. 1 (1951), 411--429.

\subsubsection{$L\log L$}
\label{subsec:ch2-l-log-l}
在命题~\ref{prop:ch2-maximal-not-l1} 的证明中, $Mf$ 的可积性在无穷远处失效, 但这并未排除局部可积性. 然而, 例子 $f(x)=x^{-1}(\log x)^{-2}\chi_{(0,1/2]}$ 表明局部可积性也可能失效. 定理~\ref{thm:ch2-local-l-log-l} 有一个部分逆命题, 它刻画了 $Mf$ 何时局部可积.

\begin{Theorem}
\label{thm:ch2-l-log-l-characterization}
若 $f$ 是支集包含于紧集 $B$ 的可积函数, 则 $Mf\in L^1(B)$ 当且仅当 $f\log^+f\in L^1(B)$.
\end{Theorem}

这一结果归功于 E. M. Stein(\emph{Note on the class $L\log L$}, Studia Math. 32 (1969), 305--310); 证明也可见 García-Cuerva 与 Rubio de Francia \MBUnresolvedReference{citation}{[6, p. 146]}. 其证明的核心是极大函数弱 $(1,1)$ 不等式的加强形式及“反向”弱 $(1,1)$ 不等式:
\begin{equation}
\label{eq:ch2-maximal-restricted-upper}
 |\{x\in\mathbb R^n:Mf(x)>\lambda\}|\le\frac C\lambda\int_{\{x:|f(x)|>\lambda/2\}}|f(x)|\,dx,
\end{equation}
\begin{equation}
\label{eq:ch2-maximal-restricted-lower}
 |\{x\in\mathbb R^n:Mf(x)>\lambda\}|\ge\frac c\lambda\int_{\{x:|f(x)|>\lambda\}}|f(x)|\,dx.
\end{equation}
不等式 \eqref{eq:ch2-maximal-restricted-upper} 由弱 $(1,1)$ 不等式应用于 $f_1=f\chi_{\{x:|f(x)|>\lambda/2\}}$ 得到; 将 $M$ 换成 $M''$ 后, 不等式 \eqref{eq:ch2-maximal-restricted-lower} 由 Calderón--Zygmund 分解得到.

在此类问题中, 希望把满足 $f\log^+f\in L^1$ 的函数视为某个 Banach 空间的元素. 但 $\int f\log^+f$ 不定义范数. 一种处理方法是引入 Luxemburg 范数: 给定 $\Omega\subset\mathbb R^n$, 定义
\begin{equation}
\label{eq:ch2-luxemburg-norm}
 \|f\|_{L\log L(\Omega)}=\inf\left\{\lambda>0:\int_\Omega\frac{|f(x)|}{\lambda}\log^+\left(\frac{|f(x)|}{\lambda}\right)dx\le1\right\}.
\end{equation}
利用这一范数可将定理~\ref{thm:ch2-l-log-l-characterization} 加强为 $\|Mf\|_{L^1(B)}\le C\|f\|_{L\log L(B)}$(见 Zygmund \MBUnresolvedReference{citation}{[21, Chapter 4]}). 以任意凸递增函数 $\Phi$ 代替 \eqref{eq:ch2-luxemburg-norm} 中的 $t\log^+t$, 得到推广 $L^p$ 空间的一类 Banach 函数空间 $L^\Phi(\Omega)$, 称为 Orlicz 空间. 这类空间具有丰富的理论; 更多资料可见 M. A. Krasnosel'skiĭ 与 Ya. B. Rutickiĭ 的著作 \emph{Convex Functions and Orlicz Spaces}, P. Noordhoff, Groningen, 1961, 以及 M. M. Rao 与 Z. D. Ren 的著作 \emph{Theory of Orlicz Spaces}, Marcel Dekker, New York, 1991.

\hypertarget{chapter-02-page-055}{}
\subsubsection{常数的大小}
\label{subsec:ch2-size-of-constants}
本章证明中, $M$ 与 $M'$ 的弱 $(1,1)$ 和强 $(p,p)$ 不等式, $1<p<\infty$, 中的常数关于维数 $n$ 呈指数增长. 对强 $(p,p)$ 不等式, 旋转法(见 \MBUnresolvedReference{corollary}{推论 4.7})给出形如 $C_pn$ 的较好常数. 进一步可证明, 存在与 $n$ 无关的常数 $C_p$, 使
\[
 \|Mf\|_p\le C_p\|f\|_p,\qquad 1<p<\infty.
\]
这一结果归功于 E. M. Stein, 证明发表于 E. M. Stein 与 J.-O. Strömberg 的论文 \emph{Behavior of maximal functions in $\mathbb R^n$ for large $n$}, Ark. Mat. 21 (1983), 259--269; 另见 E. M. Stein, \emph{Three variations on the theme of maximal functions}, 载于 \emph{Recent Progress in Fourier Analysis}, I. Peral 与 J. L. Rubio de Francia 编, pp. 229--244, North-Holland, Amsterdam, 1985. 该结果可以部分推广. 若 $B$ 是关于原点对称的凸集, 定义
\[
 M_Bf(x)=\sup_{r>0}\frac1{|rB|}\int_{rB}|f(x-y)|\,dy.
\]
则当 $p>3/2$ 时, 存在与 $B,n$ 无关的常数 $C_p$, 使
\[
 \|M_Bf\|_p\le C_p\|f\|_p.
\]
对弱 $(1,1)$ 不等式, $M$ 目前已知的最佳常数为 $n$ 阶(见上述 Stein 与 Strömberg 的论文). 对 \eqref{eq:ch2-hl-maximal-uncentered-cube} 定义的非中心极大函数(将立方体换成球), 强 $(p,p)$ 不等式中的最佳常数必随 $n$ 指数增长; 见 L. Grafakos 与 S. Montgomery-Smith, \emph{Best constants for uncentered maximal functions}, Bull. London Math. Soc. 29 (1997), 60--64.

当 $n=1$ 时, 单侧极大算子和非中心极大算子 $M''$ 的最佳常数已知, 但 $M$ 的最佳常数未知. 对 $M''$ 的弱 $(1,1)$ 不等式, 引理~\ref{lem:ch2-interval-covering} 给出 $C=2$, 而函数 $f(t)=\chi_{[0,1]}$ 表明这是最佳常数. 关于 $M''$ 的强 $(p,p)$ 不等式, 见上述 Grafakos 与 Montgomery-Smith 的论文. 关于 $M$ 的弱 $(1,1)$ 常数的上下界, 见 J. M. Aldaz, \emph{Remarks on the Hardy--Littlewood maximal function}, Proc. Roy. Soc. Edinburgh Sect. A 128 (1998), 1--9.

\hypertarget{chapter-02-page-056}{}
\subsubsection{覆盖引理}
\label{subsec:ch2-covering-lemmas}
证明 $\mathbb R^n$ 上的极大函数为弱 $(1,1)$ 型的一种标准方法是使用覆盖引理. 这里给出两个. 第一个是 N. Wiener 给出的 Vitali 型引理. 若 $B=B(x,r)$ 且 $t>0$, 记 $tB=B(x,tr)$.

\begin{Theorem}
\label{thm:ch2-wiener-covering}
设 $\{B_j\}_{j\in\mathcal J}$ 是 $\mathbb R^n$ 中的一族球. 则存在至多可数的两两不交子族 $\{B_k\}$, 使
\[
 \bigcup_{j\in\mathcal J}B_j\subset\bigcup_k5B_k.
\]
\end{Theorem}

第二个引理由 A. Besicovitch 与 A. P. Morse 独立得到; 证明及更多参考资料可见 M. de Guzmán, \emph{Differentiation of Integrals in $\mathbb R^n$}, Lecture Notes in Math. 481, Springer-Verlag, Berlin, 1985.
\begin{Theorem}
\label{thm:ch2-besicovitch-morse-covering}
设 $A$ 是 $\mathbb R^n$ 中的有界集, 且 $\{B_x\}_{x\in A}$ 是一族球, 其中 $B_x=B(x,r_x)$, $r_x>0$. 则存在至多可数的子族 $\{B_j\}$ 及只依赖维数的常数 $C_n$, 使
\[
 A\subset\bigcup_jB_j,
 \qquad
 \sum_j\chi_{B_j}(x)\le C_n.
\]
\end{Theorem}

将球换成立方体时同一结论仍成立; 更一般地, 点 $x$ 不必是球或立方体的中心, 但必须一致地靠近中心. 利用 Besicovitch--Morse 引理, 可把极大函数的结果推广到相对于其他测度的 $L^p$ 空间. 给定非负 Borel 测度 $\mu$, 定义
\[
 M_\mu f(x)=\sup_{r>0}\frac1{\mu(B_r)}\int_{B_r}|f(x-y)|\,d\mu(y).
\]
若 $\mu(B_r)=0$, 则把 $f$ 在 $B_r$ 上的 $\mu$-平均定义为 $0$. 可以证明 $M_\mu$ 相对于 $\mu$ 为弱 $(1,1)$ 型, 因而由插值知它在 $L^p(\mu)$, $1<p<\infty$, 上有界.

\hypertarget{chapter-02-page-057}{}
若以立方体定义 $M'_\mu$, 则除非 $\mu$ 满足附加倍增条件, $M_\mu$ 与 $M'_\mu$ 不必逐点等价: 存在 $C$, 使对任意球 $B$ 有 $\mu(2B)\le C\mu(B)$. 对非中心极大算子 $M''_\mu$, 当 $n>1$ 时, 除非 $\mu$ 满足倍增条件, 它甚至不必为弱 $(1,1)$ 型. 在这种情形, 弱 $(1,1)$ 不等式由 Wiener 引理推出; 当 $n=1$ 时, 即便对非倍增测度也可应用引理~\ref{lem:ch2-interval-covering}. 使 $M''_\mu$ 非弱 $(1,1)$ 型的测度例子归功于 P. Sjögren, \emph{A remark on the maximal function for measures on $\mathbb R^n$}, Amer. J. Math. 105 (1983), 1231--1233. 对某些测度, 比倍增更弱的条件也蕴含 $M''_\mu$ 为弱 $(1,1)$ 型; 例如见 A. Vargas, \emph{On the maximal function for rotation invariant measures in $\mathbb R^n$}, Studia Math. 110 (1994), 9--17.

\subsubsection{非切向 Poisson 极大函数}
\label{subsec:ch2-nontangential-poisson-maximal}
给定 $f\in L^p(\mathbb R^n)$, $u(x,t)=P_t*f(x)$ 定义 $f$ 到上半空间 $\mathbb R_+^{n+1}$ 的调和延拓. 固定 $a>0$, 令
\[
 \Gamma_a(x)=\{(y,t):|y-x|<at\}.
\]
这是顶点为 $x$、孔径为 $a$ 的锥. 我们要问是否有
\[
 \lim_{\substack{(y,t)\to(x,0)\\(y,t)\in\Gamma_a(x)}}u(y,t)=f(x)
 \quad\text{a.e. }x\in\mathbb R^n.
\]
若 $f$ 连续且具有紧支集, 则此极限对每个 $x$ 都成立. 因此, 为应用定理~\ref{thm:ch2-closed-ae-convergence-set}, 只需考虑极大函数
\[
 u_a^*(x)=\sup_{(y,t)\in\Gamma_a(x)}|u(y,t)|.
\]
存在只依赖 $a$ 的常数 $C_a$, 使
\[
 u_a^*(x)\le C_aMf(x).
\]
因此 $u_a^*$ 为弱 $(1,1)$ 型和强 $(p,p)$ 型, $1<p\le\infty$, 并可结合定理~\ref{thm:ch2-closed-ae-convergence-set} 得到几乎处处非切向收敛. 这些结果还能推广到“切向”逼近区域, 条件是相伴极大函数具有弱有界性. 例如见 A. Nagel 与 E. M. Stein, \emph{On certain maximal functions and approach regions}, Adv. in Math. 54 (1984), 83--106, 以及 D. Cruz-Uribe、C. J. Neugebauer 与 V. Olesen, \emph{Norm inequalities for the minimal and maximal operator, and differentiation of the integral}, Publ. Mat. 41 (1997), 577--604.

\subsubsection{强极大函数}
\label{subsec:ch2-strong-maximal}
令 $R(h_1,\ldots,h_n)=[-h_1,h_1]\times\cdots\times[-h_n,h_n]$. 对 $f\in L^1_{\mathrm{loc}}$, 定义
\[
 M_sf(x)=\sup_{h_1,\ldots,h_n>0}\frac1{|R(h_1,\ldots,h_n)|}\int_{R(h_1,\ldots,h_n)}|f(x-y)|\,dy.
\]

\hypertarget{chapter-02-page-058}{}
$M_s$ 在 $L^p(\mathbb R^n)$, $p>1$, 上有界, 这是单变量结果的推论. 但 $M_s$ 不是弱 $(1,1)$ 型; 最佳不等式为
\[
 |\{x\in\mathbb R^n:M_sf(x)>\lambda\}|\le C\int\frac{|f(x)|}{\lambda}\left(1+\log^+\frac{|f(x)|}{\lambda}\right)^{n-1}dx.
\]
这一结果归功于 B. Jessen、J. Marcinkiewicz 与 A. Zygmund(\emph{Note on the differentiability of multiple integrals}, Fund. Math. 25 (1935), 217--234). A. Córdoba 与 R. Fefferman 给出了几何证明(\emph{A geometric proof of the strong maximal theorem}, Ann. of Math. 102 (1975), 95--100); 另见 R. J. Bagby, \emph{Maximal functions and rearrangements: some new proofs}, Indiana Univ. Math. J. 32 (1983), 879--891. 相应的 Lebesgue 微分定理
\[
 \lim_{\max(h_i)\to0}\frac1{|R(h_1,\ldots,h_n)|}\int_{R(h_1,\ldots,h_n)}f(x-y)\,dy=f(x)\quad\text{a.e.}
\]
对某些 $f\in L^1$ 不成立, 但当 $f(1+\log^+|f|)^{n-1}$ 局部可积时成立, 特别地对某个 $p>1$ 的 $f\in L^p_{\mathrm{loc}}$ 成立.

若在 $M_s$ 的定义中允许任意方向的长方体(即平行多面体), 而不仅是边平行于坐标轴的长方体, 则所得算子在任何 $L^p$, $p<\infty$, 上都无界, 相伴微分定理甚至对有界 $f$ 也不成立. 关于这些及其他积分微分问题, 见 de Guzmán \MBUnresolvedReference{citation}{[7]}, 上述同一作者的著作, 以及 A. M. Bruckner 的专著 \emph{Differentiation of Integrals}, Amer. Math. Monthly 78 (1971), Slaught Memorial Papers, 12.

\subsubsection{Kakeya 极大函数}
\label{subsec:ch2-kakeya-maximal}
给定 $N>1$, 令 $\mathcal R_N$ 为 $\mathbb R^n$ 中所有这样的长方体所成的集合: $n-1$ 条边长为 $h$, 另一条边长为 $Nh$, $h>0$. 定义 Kakeya 极大函数
\[
 \mathcal K_Nf(x)=\sup_{x\in R\in\mathcal R_N}\frac1{|R|}\int_R|f|.
\]
由于 $\mathcal R_N$ 中每个长方体都包含于半径为 $c_nNh$ 的球中. 因而 Kakeya 极大函数受 Hardy--Littlewood 极大函数逐点控制:
\begin{equation}
\label{eq:ch2-kakeya-trivial-bound}
 \mathcal K_Nf(x)\le C_nN^{n-1}Mf(x).
\end{equation}
因此 $\mathcal K_N$ 在 $L^p$, $1<p\le\infty$, 上有界.

\hypertarget{chapter-02-page-059}{}
由 \eqref{eq:ch2-kakeya-trivial-bound} 及 $L^\infty$ 范数为 $1$, 插值得
\[
 \|\mathcal K_Nf\|_p\le C_nN^{(n-1)/p}\|f\|_p,
 \qquad 1<p<\infty.
\]
当 $p=n$ 时, 猜想该常数可改进为 $\log N$ 的幂. 由插值, 这将给出
\[
 \begin{aligned}
 p\ge n:&\quad \|\mathcal K_Nf\|_p\le C_n(\log N)^{\alpha_p}\|f\|_p,\\
 1<p\le n:&\quad \|\mathcal K_Nf\|_p\le C_nN^{n/p-1}(\log N)^{\alpha_p}\|f\|_p.
 \end{aligned}
\]
这一猜想与另外两个问题密切相关: Kakeya 集的 Hausdorff 维数, 即包含每个方向上一条线段的 Lebesgue 测度为零的集合(见 Stein \MBUnresolvedReference{citation}{[17, p. 434]}), 以及 Bochner--Riesz 乘子的有界性(见 \MBUnresolvedReference{section}{第 8 章第 5 节和第 8.3 节}).

该猜想仅在 $n=2$ 时得到完全证明; 见 A. Córdoba, \emph{The Kakeya maximal function and the spherical summation multipliers}, Amer. J. Math. 99 (1977), 1--22. 该文还讨论了它与 Bochner--Riesz 乘子的联系. Córdoba 后来证明了所有维数下 $1<p<2$ 的结果; 见 \emph{A note on Bochner--Riesz operators}, Duke Math. J. 46 (1979), 505--511. M. Christ、J. Duoandikoetxea 与 J. L. Rubio de Francia 在 \emph{Maximal operators related to the Radon transform and the Calderón--Zygmund method of rotations}, Duke Math. J. 53 (1986), 189--209 中把范围推广到 $1<p\le(n+1)/2$.

J. Bourgain 的工作极大地推动了这一问题(\emph{Besicovitch type maximal operators and applications to Fourier analysis}, Geom. Funct. Anal. 1 (1991), 147--187). 他证明了猜想对 $1<p\le(n+1)/2+\varepsilon_n$ 成立, 其中 $\varepsilon_n$ 由递推公式给出(例如 $\varepsilon_3=1/3$). T. Wolff 在 \emph{An improved bound for Kakeya type maximal functions}, Rev. Mat. Iberoamericana 11 (1995), 651--674 中把范围改进为 $1<p\le(n+2)/2$. 随后 J. Bourgain 在 \emph{On the dimension of Kakeya sets and related maximal inequalities}, Geom. Funct. Anal. 9 (1999), 256--282 中证明, 存在 $c>1/2$, 使当 $n$ 充分大时猜想对 $p\le cn$ 成立.

关于该问题近期结果及其与 Kakeya 集的联系, 可见 T. Wolff 的综述 \emph{Recent work connected with the Kakeya problem}, 载于 \emph{Prospects in Mathematics}, H. Rossi 编, pp. 129--162, Amer. Math. Soc., Providence, 1999.
